% Standalone resolution manuscript generated from solution.md.
% Compile with XeLaTeX; author and review metadata are maintained in Markdown.
\documentclass[10pt,a4paper]{article}
\usepackage[margin=24mm,headheight=15pt]{geometry}
\usepackage{fontspec}
\setmainfont{DejaVuSerif.ttf}[BoldFont=DejaVuSerif-Bold.ttf,ItalicFont=DejaVuSerif-Italic.ttf,BoldItalicFont=DejaVuSerif-BoldItalic.ttf]
\setsansfont{DejaVuSans.ttf}[BoldFont=DejaVuSans-Bold.ttf,ItalicFont=DejaVuSans-Oblique.ttf]
\setmonofont{DejaVuSansMono.ttf}
\usepackage{amsmath,amssymb,unicode-math}
\setmathfont{latinmodern-math.otf}
\usepackage{xcolor,microtype,parskip,needspace,fancyhdr}
\usepackage{bookmark}
\usepackage{xurl}
\hypersetup{colorlinks=true,urlcolor=blue,linkcolor=blue,pdftitle={A
seven-dimensional counterexample to derivative
realizability},pdfauthor={Matthew J. Colbrook},pdflang={en-GB}}
\urlstyle{same}
\setlength{\emergencystretch}{3em}
\providecommand{\tightlist}{\setlength{\itemsep}{0pt}\setlength{\parskip}{0pt}}
\providecommand{\pandocbounded}[1]{#1}
\setcounter{secnumdepth}{0}
\allowdisplaybreaks[1]
\widowpenalty=10000
\clubpenalty=10000
\pagestyle{fancy}
\fancyhf{}
\fancyhead[L]{\small\sffamily RESOLUTION / IS-03}
\fancyhead[R]{\small\sffamily 11 September 2026}
\fancyfoot[L]{\parbox[b]{0.9\textwidth}{\fontsize{8}{10}\selectfont Independent
Codex-agent verification; no external human peer review or formal
certification.}}
\fancyfoot[R]{\footnotesize\thepage}
\begin{document}
{\raggedright\Large\bfseries A seven-dimensional counterexample to
derivative realizability\par}
\vspace{0.4em}
{\normalsize\bfseries Matthew J. Colbrook\par}
{\small Department of Applied Mathematics and Theoretical Physics,
University of Cambridge, Cambridge, United Kingdom\par}
{\small\href{mailto:m.colbrook@damtp.cam.ac.uk}{m.colbrook@damtp.cam.ac.uk}\par}
\vspace{0.6em}
\textbf{Status of this manuscript:} Independently checked negative
resolution (Codex-agent review).\\
\textbf{Target:}
\href{https://github.com/ajt60gaibb/OpenProblemsInNLA/blob/b4123194697bdf6f8f82518c1dd7d6c40a30c2e0/eigenvalues-and-inverse-problems/IS-03/README.md}{IS-03,
original catalog snapshot}.\\
\textbf{Reviewed:} 11 September 2026.

The original proof draft was generated in a ChatGPT conversation. A
separate Codex agent independently checked the complete argument against
the exact target; its
\href{https://github.com/ajt60gaibb/OpenProblemsInNLA/blob/main/references/colbrook-additional-2026-09-11/verification/reviews/IS-03-review.md}{detailed
review} records \textbf{PASS}. The mathematical sections below are
unchanged from the reviewed draft. This is independent agent
verification, not external human peer review or a formal proof
certificate. The
\href{https://github.com/ajt60gaibb/OpenProblemsInNLA/blob/main/references/colbrook-additional-2026-09-11/README.md}{submission
record} preserves the original manuscript and verification history.

\section{The target and the
counterexample}\label{the-target-and-the-counterexample}

The repository asks whether, for every entrywise-nonnegative real matrix
\(A\) of order \(n\ge5\), the monic polynomial \(p_A'(z)/n\) is the
characteristic polynomial of an entrywise-nonnegative real matrix of
order \(n-1\).\footnote{\emph{Open Problems in Numerical Linear
  Algebra}, IS-03, ``Johnson's derivative-realizability conjecture,''
  checked 11 September 2026.
  \href{https://github.com/ajt60gaibb/OpenProblemsInNLA/blob/main/eigenvalues-and-inverse-problems/IS-03/README.md}{Canonical
  entry}.}

Let \(C_k\) denote the permutation matrix of a directed cycle of length
\(k\), so that \(\det(zI-C_k)=z^k-1\). Set \[
 A=\operatorname{diag}\left(\frac12,C_2,C_4\right)
 =\begin{pmatrix}
 \frac12&0&0&0&0&0&0\\
 0&0&1&0&0&0&0\\
 0&1&0&0&0&0&0\\
 0&0&0&0&1&0&0\\
 0&0&0&0&0&1&0\\
 0&0&0&0&0&0&1\\
 0&0&0&1&0&0&0
 \end{pmatrix}.
 \tag{1}
\] This is an entrywise-nonnegative real \(7\times7\) matrix. Its
characteristic polynomial is \[
 \begin{aligned}
 p(z)&=(z-\tfrac12)(z^2-1)(z^4-1)\\
 &=z^7-\tfrac12z^6-z^5+\tfrac12z^4-z^3+\tfrac12z^2+z-\tfrac12.
 \end{aligned}
 \tag{2}
\] Its normalized derivative is \[
 q(z)=\frac{p'(z)}7
 =z^6-\frac37z^5-\frac57z^4+\frac27z^3
       -\frac37z^2+\frac17z+\frac17.
 \tag{3}
\]

\textbf{Theorem 1.} There is no entrywise-nonnegative real \(6\times6\)
matrix with characteristic polynomial \(q\). Consequently, the universal
assertion in IS-03 is false.

\section{Exact negative-moment
certificate}\label{exact-negative-moment-certificate}

Let \(\mu_1,\ldots,\mu_6\) be the roots of \(q\), counted with algebraic
multiplicity, and put \(s_k=\sum_{i=1}^6\mu_i^k\). Newton's identities
determine these numbers without computing the roots.

For a monic degree-\(d\) polynomial \[
 q(z)=z^d+c_1z^{d-1}+\cdots+c_d,
\] the identities are \[
 \begin{cases}
 s_k+c_1s_{k-1}+\cdots+c_{k-1}s_1+kc_k=0,&1\le k\le d,\\
 s_k+c_1s_{k-1}+\cdots+c_ds_{k-d}=0,&k>d.
 \end{cases}
 \tag{4}
\] Applying (4) to (3) gives the following rational values: \[
 \begin{aligned}
 s_1&=\frac37,&s_2&=\frac{79}{49},&s_3&=\frac{48}{343},\\
 s_4&=\frac{6731}{2401},&s_5&=\frac{5213}{16807},&
 s_6&=\frac{219766}{117649}.
 \end{aligned}
 \tag{5}
\] The next identity is particularly short: \[
 7s_7=3s_6+5s_5-2s_4+3s_3-s_2-s_1.
 \tag{6}
\] After expressing its right-hand side over the denominator
\(117649=7^6\), its numerator is \[
 659298+182455-659638+49392-189679-50421=-8593.
\] Therefore \[
 \boxed{s_7=-\frac{8593}{823543}<0.}
 \tag{7}
\]

\textbf{Proof of Theorem 1.} If an entrywise-nonnegative matrix \(B\)
had characteristic polynomial \(q\), the trace identity for matrix
powers, valid without diagonalizability, would give \[
 \operatorname{tr}(B^7)=\sum_{i=1}^6\mu_i^7=s_7<0.
\] But every entry of \(B^7\) is nonnegative, so its trace is
nonnegative. This is a contradiction.

\section{Scope and checks}\label{scope-and-checks}

The counterexample satisfies the real field, entrywise nonnegativity,
and order condition in the repository. Neither irreducibility nor zero
trace is an assumption there. The example has positive trace and is
reducible; these are permitted, not exceptions to the target. It does
not contradict a result restricted to trace-zero matrices of orders five
or six.

The negative moment also precludes a realization after appending any
number of zero eigenvalues, since zeros do not change \(s_7\). Thus no
relaxation by zero padding repairs this particular example.

The accompanying verification script checks (2)--(7) with exact rational
arithmetic and independently computes the seventh power trace of the
companion matrix of \(q\). The companion matrix is used only to check
the certificate, not as a proposed nonnegative realization. These
diagnostics are separate from the independent proof review linked above.

The repository submission records \textbf{Solved}, with a
negative-resolution notice and the independent Codex-agent PASS report
linked above.

\section{Sources}\label{sources}

\begin{enumerate}
\def\labelenumi{\arabic{enumi}.}
\tightlist
\item
  \emph{Open Problems in Numerical Linear Algebra}, IS-03,
  \href{https://github.com/ajt60gaibb/OpenProblemsInNLA/blob/main/eigenvalues-and-inverse-problems/IS-03/README.md}{canonical
  statement}, checked 11 September 2026.
\item
  S. L. Hoover, D. McCormick, P. Paparella, and A. Thrall, \emph{On the
  realizability of the critical points of a realizable list},
  arXiv:1712.05454, Conjecture 1.2 and §6; \emph{Linear Algebra and its
  Applications} 555 (2018), 301--313.
  \href{https://arxiv.org/abs/1712.05454}{Manuscript};
  \href{https://doi.org/10.1016/j.laa.2018.06.024}{publication}. This is
  the primary reference identified by the repository for the conjecture,
  not a source for the counterexample above.
\end{enumerate}
\end{document}
